\section{OBJECTIVE AND SOURCE REVIEW}
\textbf{Erd\H{o}s \#623; independent attempt, 2026-09-23.}
For $|X|=\aleph_\omega$ and $f:[X]^{<\omega}\to X$ with $f(A)\notin A$, must there be an infinite $Y$ such that $f(B)\notin Y$ for every finite $B\subseteq Y$? I read the complete \texttt{623.tex}. Its finite free-set argument does not automatically produce a compatible infinite sequence. I give an explicit smaller-cardinality model demonstrating the exact failure of that inference.
\section{WORK: FINITE FREEDOM WITHOUT INFINITE FREEDOM}
On $X=\mathbb N_0$, define
\[
f(A)=\begin{cases}|A|,&|A|\notin A,\\
\min(\mathbb N_0\setminus A),&|A|\in A.
\end{cases}
\]
This always lies outside $A$. For any fixed $r$, take $Y_r=\{r+1,\ldots,2r\}$. Every $B\subseteq Y_r$ has $|B|\le r$, so $|B|\notin B$ and $f(B)=|B|\notin Y_r$. Thus independent sets of every finite size exist.
But no infinite independent $Y$ exists. Choose $m\in Y$ and choose $A\subseteq Y\setminus\{m\}$ of size $m$, possible by infinitude (including $A=\varnothing$ if $m=0$). Then $f(A)=m\in Y$, a contradiction.
\section{ADVERSARIAL VERIFICATION AND REMAINING GAP}
This is a counterexample on a countable set only, not on $\aleph_\omega$. Extending it by assigning a numerical colour to each point does not preserve the argument, because $f(A)$ must be an actual point of the proposed infinite set, not merely share its colour. The tree of finite independent sets has arbitrarily high levels but can branch infinitely, so the finite-branching form of K\"onig's lemma does not apply. The source's finite Kuratowski theorem remains compatible with this example. A construction or theorem using the specific singular cardinal $\aleph_\omega$ is still required.
\section{SOURCES CHECKED}
Complete local source \texttt{gpt\_pro\_5.4/623.tex}; exact statement checked in \url{https://raw.githubusercontent.com/google-deepmind/formal-conjectures/main/FormalConjectures/ErdosProblems/623.lean}. The explicit countable example is proved directly and no literature novelty is claimed.
\section{FINAL}
\textbf{UNRESOLVED.} The invalid compactness inference is concretely blocked; the cardinality-specific question remains.
\section{COMPLETION ESTIMATE}
Conservative subjective progress toward the stated uncountable problem.
\textbf{COMPLETION: 6\%}
